% =====================================================================
% 5. Learning strategies and generator: implementation (MLLB, Maria Martinez)
% Shared T13, T14, T24 (DECISIONS D20, block note 08); generator (Oriana,
% T23, D19, block note 13), summarised and attributed.
% =====================================================================
\section{Learning strategies and generator: implementation}\label{sec:mllb-strategies}

\subsection{Transfer from the self-supervised generator}

Oriana and I ran the transfer experiments together. The pre-trained network is Oriana's ECG-GPT
(Section~\ref{sec:mllb-generator}); its \texttt{record\_features} method returns one vector of 128 numbers
per record, the mean of the outputs of its final LayerNorm over the windows that cover lead II. Listing~\ref{lst:transfer}
puts a new head on top.

\lstinputlisting[firstnumber=185,caption={Generator backbone with a new classification head (\texttt{src/transfer.py}).},label={lst:transfer}]{common/code/transfer_classifier.py}

\begin{itemize}[leftmargin=1.2em,itemsep=1pt]
  \item \textbf{Freezing.} For the linear probe every backbone parameter gets \texttt{requires\_grad\_(False)};
        since the optimiser only receives trainable parameters, the backbone cannot change.
        Freezing the weights is not enough, though: the loop calls \texttt{model.train()} every epoch, which
        would switch on the backbone's dropout and give a different vector for the same record each time.
        The class therefore overrides \texttt{train()} to keep the frozen backbone in eval mode.
  \item \textbf{Caching the features.} With a frozen backbone each record always gives the same vector, so
        the vectors are computed once per checkpoint and the head is trained on them; a test checks that this
        is the same model as passing the tokens every epoch.
  \item \textbf{Fine-tuning.} Everything trainable, with a learning rate of $10^{-4}$ instead of $10^{-3}$, as in
        the class transfer-learning notebook \cite{course_nb_transfer}: at the start the new head is random and
        its large gradients flow into the pre-trained layers. The ablation at $10^{-3}$ showed that this was
        not essential here.
  \item \textbf{Memory.} One record is four windows of 250 tokens, so the activations kept for
        backpropagation grow fast; the transfer runs use batches of 32 records instead of 64, the same for
        every transfer configuration.
  \item \textbf{The 10\,\% subset} is drawn by patient with a generator seeded by the run's seed, and is the
        same for the three strategies of that seed, so they are compared on the same labelled records.
\end{itemize}
Results (Table~\ref{tab:transfer_main}): fine-tuning reached
\transferTestAurocMeanFinetuneFOneZeroZero{} on test with all the labels and
\transferTestAurocMeanFinetuneFOneZero{} with 10\,\%, against \transferTestAurocMeanScratchFOneZeroZero{}
and \transferTestAurocMeanScratchFOneZero{} from scratch.

% Shared float (common/make_floats.py). Table: tables/transfer_main.tex
\begin{table}[htbp]
\centering
\caption[Transfer from the ECG-GPT: main results]{ECG-GPT reused for the five superclasses (lead II): the three strategies with 100\,\% and 10\,\% of the training labels (mean $\pm$ SD over seeds 42--44).}
\label{tab:transfer_main}
\small
\begin{adjustbox}{max width=\textwidth}
% Generated by src/transfer_experiments.py from results/logs/transfer_summary.json.
% Needs report/numbers.tex (alias "transfer") and booktabs.
\begin{tabular}{llrrccc}
\toprule
Labels & Strategy & Train.\ param. & Train records & Val.\ macro AUROC & Test macro AUROC & Test macro F1 \\
\midrule
100\% & From scratch & \transferParamsScratchFOneZeroZero{} & \transferNTrainScratchFOneZeroZero{} & \transferValAurocMeanScratchFOneZeroZero{} $\pm$ \transferValAurocSdScratchFOneZeroZero{} & \transferTestAurocMeanScratchFOneZeroZero{} $\pm$ \transferTestAurocSdScratchFOneZeroZero{} & \transferTestFOneMeanScratchFOneZeroZero{} $\pm$ \transferTestFOneSdScratchFOneZeroZero{} \\
100\% & Linear probe (frozen) & \transferParamsProbeFOneZeroZero{} & \transferNTrainProbeFOneZeroZero{} & \transferValAurocMeanProbeFOneZeroZero{} $\pm$ \transferValAurocSdProbeFOneZeroZero{} & \transferTestAurocMeanProbeFOneZeroZero{} $\pm$ \transferTestAurocSdProbeFOneZeroZero{} & \transferTestFOneMeanProbeFOneZeroZero{} $\pm$ \transferTestFOneSdProbeFOneZeroZero{} \\
100\% & Fine-tuning & \transferParamsFinetuneFOneZeroZero{} & \transferNTrainFinetuneFOneZeroZero{} & \transferValAurocMeanFinetuneFOneZeroZero{} $\pm$ \transferValAurocSdFinetuneFOneZeroZero{} & \transferTestAurocMeanFinetuneFOneZeroZero{} $\pm$ \transferTestAurocSdFinetuneFOneZeroZero{} & \transferTestFOneMeanFinetuneFOneZeroZero{} $\pm$ \transferTestFOneSdFinetuneFOneZeroZero{} \\
\addlinespace
10\% & From scratch & \transferParamsScratchFOneZero{} & \transferNTrainScratchFOneZero{} & \transferValAurocMeanScratchFOneZero{} $\pm$ \transferValAurocSdScratchFOneZero{} & \transferTestAurocMeanScratchFOneZero{} $\pm$ \transferTestAurocSdScratchFOneZero{} & \transferTestFOneMeanScratchFOneZero{} $\pm$ \transferTestFOneSdScratchFOneZero{} \\
10\% & Linear probe (frozen) & \transferParamsProbeFOneZero{} & \transferNTrainProbeFOneZero{} & \transferValAurocMeanProbeFOneZero{} $\pm$ \transferValAurocSdProbeFOneZero{} & \transferTestAurocMeanProbeFOneZero{} $\pm$ \transferTestAurocSdProbeFOneZero{} & \transferTestFOneMeanProbeFOneZero{} $\pm$ \transferTestFOneSdProbeFOneZero{} \\
10\% & Fine-tuning & \transferParamsFinetuneFOneZero{} & \transferNTrainFinetuneFOneZero{} & \transferValAurocMeanFinetuneFOneZero{} $\pm$ \transferValAurocSdFinetuneFOneZero{} & \transferTestAurocMeanFinetuneFOneZero{} $\pm$ \transferTestAurocSdFinetuneFOneZero{} & \transferTestFOneMeanFinetuneFOneZero{} $\pm$ \transferTestFOneSdFinetuneFOneZero{} \\
\bottomrule
\end{tabular}

\end{adjustbox}
\end{table}


\subsection{Single lead and zero-shot}

The single-lead ablation needed no new model code: every architecture takes \texttt{n\_leads} as an
argument, so the final model of each family was rebuilt with \texttt{n\_leads=1} and its 12-lead arguments,
and fed lead II only. Runs are named \texttt{onelead\_*} so that they can never be mixed with the 12-lead
runs that each block's summary reads.

The zero-shot experiment loads CLIP ViT-B/32 with OpenAI's weights through \texttt{open\_clip}, as
\texttt{ViT-B-32-quickgelu}: OpenAI trained this model with the QuickGELU activation, and the plain
\texttt{ViT-B-32} configuration of \texttt{open\_clip} uses GELU, so it would not have been exactly the model
of the class notebook \cite{course_nb_clip}. Each ECG is drawn as a $224 \times 224$ image, the size CLIP
expects, so its preprocessing does not crop it. The first run failed because the image features had been
moved to the CPU and the text features stayed on the GPU; the fix is described in
Appendix~\ref{app:ai-code}.

\subsection{The generator (Oriana's block)}\label{sec:mllb-generator}

The generator was implemented by Oriana; it reuses my attention module with \texttt{causal=True}. Lead II is
quantised on a microvolt grid into integer tokens; the training windows of 250 tokens are cut at random
positions inside \texttt{\_\_getitem\_\_}, with a NumPy generator seeded by the run, so the draws are
reproducible with \texttt{num\_workers=0}. The model returns logits of shape (batch, $T$, $V$), which must
be flattened to (batch $\cdot$ $T$, $V$) before \texttt{cross\_entropy}, otherwise PyTorch reads $T$ as the
class axis. Generation runs under \texttt{torch.no\_grad()} in eval mode, samples the next token with
\texttt{torch.multinomial} and a seeded generator, and only computes the output layer at the last position.
Its main implementation lesson concerns memory: the attention matrix grows with the square of the
context, and generating with long contexts ran out of the 4\,GB of the GPU until the batch was reduced in
proportion.
